% section: 線形漸化式と線形代数

\section{線形漸化式と線形代数}

\subsection{解の集合には構造がある}

二階線形漸化式
\[
  a_{n+2}=pa_{n+1}+qa_n
\]
を考える。
二つの解 $a_n,b_n$ があれば、定数 $c,d$ に対して
\[
  c a_n+d b_n
\]
も解になる。
これを線形性という。

さらに、初期値 $a_0,a_1$ を指定すれば、以後の項は一意に決まる。
したがって、解の全体は「二つの初期値を選ぶ空間」と同じくらいの自由度を持つ。
このことを、解の空間の次元が $2$ である、と表現する。

特性方程式の二つの解から一般項を作る方法は、単に計算がうまくいくからではない。
二次の漸化式の解全体を、二つの基本的な解の組み合わせで表しているのである。

\subsection{特性方程式と固有値}

\[
  \boldsymbol v_n=
  \begin{pmatrix}a_{n+1}\\a_n\end{pmatrix}
\]
と置けば、
\[
  \boldsymbol v_{n+1}
  =
  \begin{pmatrix}p&q\\1&0\end{pmatrix}
  \boldsymbol v_n
\]
となる。
この行列を $A$ とする。

もし
\[
  A\boldsymbol v=\lambda\boldsymbol v
\]
を満たすベクトルがあれば、その方向の運動は
\[
  \boldsymbol v,\ \lambda\boldsymbol v,\ \lambda^2\boldsymbol v,\ldots
\]
となる。
行列による更新が、数の等比的な更新に変わった。

固有値を求める条件は、
\[
  \det(A-\lambda I)=0
\]
であり、
\[
  \lambda^2-p\lambda-q=0
\]
となる。
これは二階漸化式の特性方程式と一致する。
特性方程式は、行列の固有値を求める方程式としても現れているのである。

\subsection{重解のときに $n$ が現れる理由}

特性方程式が重解 $\lambda$ を持つ場合、二つの異なる等比数列だけでは解の自由度を埋められない。
例えば、
\[
  a_{n+2}-2a_{n+1}+a_n=0
\]
では、差分を取ると、
\[
  \Delta^2a_n=0
\]
である。
従って、
\[
  \Delta a_n=B
\]
は定数であり、
\[
  a_n=A+Bn
\]
となる。

特性方程式は
\[
  (\lambda-1)^2=0
\]
である。
重解 $1$ に対して、$1^n$ だけでなく $n1^n$ が現れた。
この $n$ は、差分を一回戻すときに現れる一次的な増加を記録している。

一般に、特性根 $r$ が $m$ 重に現れると、
\[
  r^n,\ nr^n,\ n^2r^n,\ldots,n^{m-1}r^n
\]
が解を作る。
大学で学ぶジョルダン標準形は、この現象を行列の側から整理したものである。

\subsection{複素数の根と振動}

特性根が
\[
  r=\rho(\cos\theta+i\sin\theta)
\]
と
\[
  \overline r=\rho(\cos\theta-i\sin\theta)
\]
であるとする。
オイラーの公式
\[
  e^{i\theta}=\cos\theta+i\sin\theta
\]
を使えば、二つの根は $\rho e^{i\theta}$ と $\rho e^{-i\theta}$ である。

このとき、実数解は
\[
  a_n=\rho^n\{C\cos(n\theta)+D\sin(n\theta)\}
\]
という形になる。
$\rho^n$ は振幅の増減を、$\cos(n\theta)$ と $\sin(n\theta)$ は振動を表す。

実根だけを考えていたときには見えなかった振動が、複素根の偏角として現れた。
複素数は、実数の数列に存在する振動を記録するための自然な言語になっている。
